\section{The \texorpdfstring{$\varepsilon$}{epsilon} Primitive}
\label{sec:epsilon}

\subsection{Token-Level Uncertainty}
\label{sec:eps-token}

Modern LLM APIs return, alongside each generated token, a log-probability
distribution over the top-$k$ alternative tokens at that position. For a
token at position $t$ with top-$k$ alternatives $\{(v_1, p_1), \ldots, (v_k, p_k)\}$:

\begin{equation}
  \eps_t = \frac{-\sum_{i=1}^{k} p_i \log p_i}{\log k}
  \label{eq:eps-token}
\end{equation}

This is the normalized partial Shannon entropy over the top-$k$ distribution.
$\eps_t \in [0, 1]$: zero when the model is certain (all probability mass on
one token); one when probability is uniform across all $k$ alternatives.
We use $k = 5$ throughout, matching the maximum provided by the OpenAI API.

\paragraph{Why partial entropy.} The full vocabulary entropy is unavailable
from the API --- only the top-$k$ log-probabilities are returned. Partial
entropy over the top-5 is an underestimate of full entropy, but it captures
the consequential signal: if the model is split between two APIs, those two
alternatives appear in the top-5 with high probability mass. A model that is
certain produces $\eps_t \approx 0$ even under partial measurement.

\paragraph{Connection to the name.} The use of $\eps$ (epsilon) is deliberate.
In numerical analysis, epsilon denotes a small perturbation that propagates
through a computation. A high-$\eps$ token is a perturbation at a specific
decision point: the model was nearly as likely to choose differently, and that
different choice would have propagated through the rest of the generation.

\subsection{Three Types of Code Generation Uncertainty}
\label{sec:eps-taxonomy}

Empirical analysis of our 30-prompt calibration benchmark reveals three
categorically distinct types of token-level uncertainty in generated code.
This taxonomy is the foundation for the filtering approach in
Section~\ref{sec:eps-filtering}.

\paragraph{Type 1 --- Cosmetic uncertainty (naming).}
The model is uncertain about what to \emph{call} something. Example: splitting
between \texttt{sort\_list}, \texttt{sort\_integers}, and \texttt{sort\_numbers}
for a function that sorts integers. The token-level $\eps$ is real and can be
high (0.7--0.9), but any choice produces equally correct code. This uncertainty
is driven by the model having no naming convention to follow. It dominates
naive entropy measurements on simple utility function prompts: in our benchmark,
the top uncertainty flag for 9 of 10 LOW prompts was a function or parameter
name on the first line of the generated function.

\paragraph{Type 2 --- Consequential uncertainty (API and library decisions).}
The model is uncertain which \emph{API or library} to use, where the choice
affects correctness, compatibility, or safety. Example: splitting between
\texttt{stripe.Charge.create()} ($\eps = 0.73$, probability 0.52) and
\texttt{stripe.PaymentIntent.create()} (probability 0.47). This is the
target signal. The split reflects genuine uncertainty about which version of
the API is appropriate --- uncertainty that no naming convention resolves,
because the split is semantic rather than stylistic.

\paragraph{Type 3 --- Implementation-approach uncertainty (algorithm structure).}
The model is uncertain which \emph{algorithm or structure} to use, where
multiple approaches are all correct. Example: recursive vs.\ iterative
Fibonacci; for-loop vs.\ list comprehension for deduplication. The token-level
split is real, and the choices produce structurally different code, but all
options produce correct output for well-specified problems. This type is
irreducible without semantic knowledge of the function specification. It
constitutes the residual false positive floor on simple utility function prompts.

The practical significance of this taxonomy is that only Type~2 uncertainty
correlates with real risk. Types~1 and~3 are noise to be minimized.

\subsection{AST-Based Cosmetic Filtering}
\label{sec:eps-filtering}

We eliminate Type~1 uncertainty using the generated code's abstract syntax
tree. After stripping markdown fences and parsing with Python's \texttt{ast}
module, we identify \emph{declaration positions}: source ranges where the
model is choosing a name rather than an API.

\textbf{Filtered (cosmetic):}
\begin{itemize}[noitemsep]
  \item Function and method names (the identifier following \texttt{def})
  \item Parameter names (\texttt{ast.arg} nodes)
  \item Local variable names (\texttt{ast.Name} in \texttt{Store} context:
    assignments, \texttt{for} targets, \texttt{with}-\texttt{as} targets,
    comprehension targets)
\end{itemize}

\textbf{Not filtered (consequential):}
\begin{itemize}[noitemsep]
  \item Import targets (\texttt{ast.alias}) --- module and symbol selection matters
  \item Attribute access (\texttt{ast.Attribute}) --- method and property selection matters
  \item Keywords (\texttt{async}, \texttt{def}, \texttt{return}, \ldots) --- structural choices matter
  \item Type annotations --- may carry compatibility information
\end{itemize}

Token-to-declaration matching uses range overlap rather than exact column
matching, because BPE tokenization introduces a leading space (the token
\texttt{~find} starts at column 3 when the identifier \texttt{find\_max}
starts at column 4). Multi-token identifiers (\texttt{find} $+$ \texttt{\_max})
are handled naturally by range overlap.

A pure-identifier guard prevents over-filtering: tokens such as \texttt{(db}
combine a structurally-significant delimiter (\texttt{(}) with the first
character of a parameter name. Only tokens where the stripped value is a
pure identifier (or a delimiter-prefixed identifier where all top alternatives
share the same delimiter) are excluded.

\textbf{Effect on benchmark.} Table~\ref{tab:filtering} shows the effect of
each filtering stage on the LOW-prompt false positive rate.

\begin{table}[h]
\centering
\caption{False positive rate on LOW prompts (10 simple utility functions)
         under successive filtering stages.}
\label{tab:filtering}
\begin{tabular}{lc}
\toprule
\textbf{Configuration} & \textbf{FP Rate} \\
\midrule
Naive token entropy (no filtering)      & 90\% \\
+ Whitespace / comment filtering        & 90\% \\
+ System prompt naming conventions      & 60\% \\
+ AST declaration filtering             & 60\% \\
+ Fence-format stripping                & 50\% \\
\midrule
Residual (Type 3 only)                  & 50\% \\
\bottomrule
\end{tabular}
\end{table}

The naming convention and AST filter address the same underlying phenomenon
--- cosmetic token uncertainty --- through complementary mechanisms: the
system prompt suppresses uncertain naming at generation time, while AST
filtering removes cosmetic tokens that survive generation. Their effects are
largely overlapping on this benchmark. When the system prompt naming
convention is removed, the AST filter alone reduces the FP rate from 90\% to
70\%; the two approaches combined reduce it to 50\%. AST filtering's
independent contribution is most visible in adversarial settings where the
system prompt is absent or ignored by the model.

The residual 50\% consists entirely of Type~3 implementation-approach
uncertainty on simple utility function prompts (Fibonacci, list operations,
binary search). These prompts are not representative of production code
generation tasks (API integration, library use) where Type~3 uncertainty
is substantially rarer.

\subsection{File- and Function-Level Aggregation}
\label{sec:eps-aggregation}

\paragraph{Max aggregation.}
The file- or function-level $\eps$ is the maximum token-level $\eps$ over
all consequential code tokens above a noise floor $\delta$:

\begin{equation}
  \eps_{\text{file}} = \max \{ \eps_t : t \text{ is a code token} \wedge \eps_t > \delta \}
  \label{eq:eps-file}
\end{equation}

We use $\delta = 0.30$ throughout. Tokens with $\eps_t \leq \delta$ represent
decisions where the model's top alternative had probability below $e^{-0.30} \approx 0.74$
of the chosen token --- effectively certain choices.

The compound alternative $1 - \prod(1 - \eps_t)$ was evaluated and rejected:
for a 200-token function with even moderate per-token entropy, the compound
score converges to $1.0$ regardless of whether any single decision is
genuinely risky. Max is bounded, interpretable, and answers the relevant question for
code review: \emph{how uncertain was the model at its single worst
consequential decision?}

\paragraph{Per-function aggregation.}
Function boundaries are detected from \texttt{def} and \texttt{async def}
statements in the generated code. For multi-function outputs, $\eps$ is
computed independently per function, enabling localization at function
granularity.

\paragraph{Max vs.\ mean ablation.}
An alternative aggregation --- mean $\eps$ over all consequential tokens above
$\delta$ --- was evaluated empirically across all 90 runs of the 30-prompt
benchmark (3 models $\times$ 30 prompts). Max and mean produced
\emph{identical} fire/no-fire decisions for every prompt on every model.
When $\eps$ fires, the peak token dominates the mean substantially (e.g., a
single token with $\eps_t = 0.908$ in a 5-token contributing set yields
$\eps_\text{mean} \approx 0.60$, well above any plausible threshold). When
$\eps_\text{file} = 0.000$, no token contributes and both measures are zero.
This empirical equivalence validates the max-based design: max preserves the
interpretability of pointing to the single worst consequential decision while
producing the same practical outcomes as mean.

\subsection{Status Thresholds and Intervention Protocol}
\label{sec:eps-thresholds}

\begin{table}[h]
\centering
\caption{$\varepsilon$ status thresholds and recommended developer actions.
         Thresholds calibrated empirically against GPT-4o-2024-08-06.}
\label{tab:thresholds}
\begin{tabular}{llp{5.5cm}}
\toprule
\textbf{Status} & \textbf{Condition} & \textbf{Recommended action} \\
\midrule
\textsc{complete} & $\eps < 0.30$ & No intervention required \\
\textsc{flagged}  & $0.30 \leq \eps < 0.65$ & Developer reviews flagged tokens \\
\textsc{paused}   & $0.65 \leq \eps < 0.95$ & Pause workflow; present alternatives \\
\textsc{aborted}  & $\eps \geq 0.95$ & Refuse output; re-prompt or escalate \\
\bottomrule
\end{tabular}
\end{table}

\paragraph{Adaptive ensemble threshold.}
After sufficient historical runs on similar prompts (retrieved via
$k$-nearest-neighbor lookup in embedding space), $\eps$ transitions from
absolute thresholds to a statistically-derived threshold. The transition uses
a two-condition domain switch:

\begin{itemize}[noitemsep]
  \item \textbf{MAD domain} ($n \geq 5$ neighbors, $\geq 500$ cumulative code tokens):
    threshold $= \text{median}(\eps) + 3.0 \cdot \text{MAD}(\eps) / 0.6745$,
    the modified Z-score of Iglewicz and Hoaglin~\cite{iglewicz1993detect}.
  \item \textbf{Conformal domain} ($n \geq 20$ neighbors, $\geq 3{,}000$ cumulative
    code tokens): empirical 95th-percentile of the neighborhood $\eps$
    distribution, providing a conformal coverage guarantee.
\end{itemize}

Calibration on 16 Scenario~E runs shows that the MAD domain activates at
run 13 (totaling 515 code tokens across the neighborhood), as predicted by
the token-budget formula. For this scenario, where $\eps_\text{file} \in
[0.69, 0.92]$ on every run, the MAD threshold stabilizes above the absolute
hard threshold (0.65), providing no additional discrimination --- correctly,
since the absolute threshold already fires on every run. The adaptive
ensemble adds value in heterogeneous prompt histories where some prompt
classes produce high-$\eps$ generations and others do not; in such settings
it prevents the absolute threshold from firing on structurally low-$\eps$
domains that happen to produce the occasional high-entropy token.
Scenario~E, being uniformly high-uncertainty, saturates the adaptive
threshold immediately.
